\documentclass[a4paper,12pt]{article}
\usepackage[UTF8]{ctex}
\usepackage{geometry}
\geometry{left=2.5cm,right=2.5cm,top=2.5cm,bottom=2.5cm}
\usepackage{graphicx}
\usepackage{caption}
\usepackage{booktabs}
\usepackage{tabularx}
\usepackage{float}
\usepackage{multirow}
\usepackage{array}
\usepackage{xcolor}

\begin{document}

% 附件2 封面
\begin{titlepage}
    \vspace*{1cm}
    \noindent 附件 2 \hfill \Large 序号：\underline{\makebox[3cm]{2024GDJSDN}}\\[2.5cm]
    
    \begin{center}
        {\Huge \textbf{第十二届全国大学生光电设计竞赛}}\\[0.8cm]
        {\Huge \textbf{创意组}}\\[2.5cm]
        
        {\fontsize{36pt}{40pt}\selectfont \textbf{创意策划书}}\\[3cm]
        
        {\Large \textbf{竞赛题目：}\underline{\makebox[12cm]{“让紫蛋飞 ”——基于紫外分光法的尿液中蛋白质检测}}}\\[1cm]
        {\Large \textbf{竞赛队名：}\underline{\makebox[12cm]{紫色分旦}}}\\[2cm]
    \end{center}
    
    \begin{table}[H]
        \centering
        \renewcommand{\arraystretch}{1.8}
        \begin{tabular}{|p{15cm}|}
        \hline
        \textbf{指导教师推荐理由：} \\
        \hline
        1、通过紫外分光检测尿蛋白浓度，具有较高的精确度 \\
        \hline
        2、检测方法操作简便，可即时获得检测结果 \\
        \hline
        3、相较市面同类产品，具有成本优势和较高商业价值及市场前景 \\
        \hline
        \end{tabular}
    \end{table}
    
    \vfill
    \begin{center}
        {\Large \textbf{全国大学生光电设计竞赛委员会制}}\\[0.5cm]
        {\Large \textbf{二〇二四年五月}}
    \end{center}
\end{titlepage}

\newpage

% 成员表格
\begin{table}[H]
    \centering
    \renewcommand{\arraystretch}{1.5}
    \begin{tabular}{|c|c|c|c|c|c|c|}
        \hline
        \multicolumn{2}{|c|}{队长姓名} & 孙以清 & 学校 & 复旦大学 & 学号 & \\
        \hline
        \multicolumn{2}{|c|}{专业班级} & \multicolumn{3}{c|}{电子信息科学与技术} & 手机 & \\
        \hline
        \multicolumn{2}{|c|}{E-mail} & \multicolumn{5}{c|}{} \\
        \hline
        \multirow{5}{*}{合作者} & 姓名 & 学号 & 专业、班级 & \multicolumn{2}{c|}{所在学校} & 联系方式 \\
        \cline{2-7}
        & 周润 & & 电子信息科学与技术 & \multicolumn{2}{c|}{复旦大学} & \\
        \cline{2-7}
        & 刘志杰 & & 电子信息科学与技术 & \multicolumn{2}{c|}{复旦大学} & \\
        \cline{2-7}
        & & & & \multicolumn{2}{c|}{} & \\
        \cline{2-7}
        & & & & \multicolumn{2}{c|}{} & \\
        \hline
        \multirow{2}{*}{指导教师} & 姓名 & 学校 & 所在院系 & \multicolumn{2}{c|}{联系电话/手机} & E-mail \\
        \cline{2-7}
        & 倪刚 & 复旦大学 & 信息科学与工程学院 光科学与工程系 & \multicolumn{2}{c|}{} & \\
        \hline
    \end{tabular}
\end{table}

\vspace{1cm}

\section*{一、 产品介绍}
近年来，随着我国人口老龄化加剧，国民慢性肾病患病率逐渐上升，然而多数患病者对自己的病情知之甚少，因而对于便捷、快速、经济的尿液检测设备的市场需求正在上升。本团队拟设计一款低成本、可反复利用、便携式的尿液中蛋白质检测设备，用于评估潜在患者的肾脏病情况，可广泛应用于家庭健康管理、养老院和社区医疗服务等领域。

本便携式尿液中蛋白质检测设备利用紫外分光光度法，通过 280-285nm 波长的紫外光透射率测定尿液中的蛋白质浓度。该设备操作简便、成本低廉且可重复使用，旨在提供快速、经济、高效的尿液蛋白质含量检测方案。

本设备主要由以下部分组成：紫外 LED、菲涅尔透镜、比色皿、光电二极管及电路、ESP32 单片机、显示屏。其中通过紫外 LED 和菲涅尔透镜产生聚焦的紫外光，透过比色皿中的尿液，由光电二极管接收透过比色皿的光信号，经过数据处理后反馈至单片机，最终将结果即时显示在显示屏上。

\begin{figure}[H]
    \centering
    \fbox{\begin{minipage}{0.9\textwidth}
    \centering
    \vspace{0.5cm}
    \textcolor{gray}{\LARGE [此处放置图片: 产品实物示意图及显示屏示意图]}\\[0.5cm]
    \textbf{【图片内容】：} 包含两个部分：1. 产品实物俯视图，标注了紫外LED、菲涅尔透镜、比色皿、光电二极管及电路、ESP32单片机的位置。 2. 显示屏侧视图。\\
    \textbf{【图片作用】：} 直观展示了检测设备的内部硬件组成架构、各元件的相对位置以及最终的实物雏形，帮助理解产品的光路设计和物理结构。\\
    \vspace{0.5cm}
    \end{minipage}}
    \caption{产品实物示意图（俯视图）及显示屏示意图（侧视图）}
\end{figure}

\section*{二、 市场调研}
\subsection*{2.1 需求分析}
\textbf{（1）老龄化社会与健康监测需求}\\
随着我国人口老龄化的不断加剧，65岁以上老年人口数量正在逐年增加，预计未来10年65岁以上老龄人口可增长至3亿人，占总人口比重持续攀升。因此，对于老年群体的健康管理逐渐成为社会关注的重点。同时，慢性肾脏病(CKD)在老年人中的发病率显著高于年轻人群，以北京地区为例，60～69岁、70～79岁和80岁以上老年人群CKD的患病率分别为20.8\%、30.5\%和37.8\%。而蛋白尿作为老年CKD的主要危险因素，将会导致病患心力衰竭、卒中、新发冠心病和总心血管事件的发生率明显提高。因而尿液中蛋白质含量的检测需求变得愈发重要。传统检测方法往往费时费力且成本高昂，因此市场急需一种快速、简便、经济的检测设备。

\textbf{（2）医疗养老机构与家庭的需求}\\
随着老龄化社会的发展，养老院和社区医疗机构的数量不断增加。政府大力支持养老产业的发展，推动医养结合、智慧养老等改革，而这些政策为各种各样的健康监测设备带来了庞大的市场需求。同样地，近年来，各个家庭健康管理意识有了显著提升，这也使得便携式健康检测设备市场需求增加。家庭用户需要一种方便上手、价格适中的检测设备，用于家庭成员的日常健康监测。本便携式尿液蛋白质检测设备正好契合这些需求，能够在医疗养老机构以及家庭环境中方便快捷地进行尿液检测，及时发现健康问题。

\subsection*{2.2 市场分析}
\textbf{（1）医疗设备市场}\\
据统计，中国医疗设备市场规模不断扩大，2022年达到9.4万亿元。尿液检测设备作为医疗设备的重要组成部分，其市场前景广阔。随着分级诊疗制度的推进，基层医疗机构对便携式尿液检测设备的需求也在增加。预计未来几年内，尿液检测设备市场将继续保持增长态势。

\textbf{（2）家庭健康监测市场}\\
随着健康意识的普及和家庭健康管理需求的增加，家庭健康监测设备市场也在快速增长。便携式尿液蛋白质检测设备的出现，将进一步推动这一市场的发展。根据市场研究数据，家庭健康监测设备市场的复合年增长率预计将保持在10\%以上，市场潜力巨大。

\section*{三、 产品技术}
\subsection*{3.1 国内外技术与产品分析}
目前，液体中蛋白质含量的检测技术主要有紫外吸收法、微量凯氏定氮法、双缩脲法、Lowry法、考马斯亮蓝法、HPLC法等。其中，微量凯氏定氮法需要较长的实验时间，且灵敏度较低；双缩脲法与Lowry法同样具有耗时长的特点，并且需要精密的操作；考马斯亮蓝法虽然灵敏度高，但不适用于多种蛋白质的检测；HPLC法有着定量准确、灵敏度高、重复性好、自动化程度高的优点，但需要同种蛋白质作为标准品，因此多用于成品质量控制。

在医疗保健行业中常用的设备是尿液分析仪，其测试蛋白质含量的原理基于指示剂蛋白误差原理，膜块中主要含有酸碱性指示剂如溴酚蓝、枸橼酸缓冲系统和表面活性剂。在pH值为3.2时，溴酚蓝产生的阴离子与带阳离子的蛋白质（如白蛋白）结合，导致颜色变化。再通过光的吸收和反射，使用比色法进行检测。但是这种方法存在着设备昂贵的缺点。

\begin{table}[H]
    \centering
    \caption{检测方法或设备及其特点}
    \renewcommand{\arraystretch}{1.2}
    \begin{tabular}{|p{4cm}|p{10cm}|}
    \hline
    \textbf{检测方法或设备} & \textbf{方法特点} \\
    \hline
    微量凯氏定氮法 & 较长的实验时间，且灵敏度较低 \\
    \hline
    双缩脲法与 Lowry 法 & 具有耗时长的特点，并且需要精密的操作 \\
    \hline
    考马斯亮蓝法、HPLC 法 & 灵敏度高，但不适用于多种蛋白质的检测 \\
    \hline
    尿液分析仪 & 针对多项指标，原理复杂，设备昂贵 \\
    \hline
    紫外光吸收法 & 灵敏，快速、易受干扰 \\
    \hline
    \end{tabular}
\end{table}

\subsection*{3.2 技术方案论证}
本设备通过紫外光吸收法进行蛋白质含量检测。紫外吸收法是通过利用蛋白质分子中的芳香族氨基酸（如色氨酸、赖氨酸、苯丙氨酸）在280nm波长处的最大吸收峰。通过分光光度的方法检测280nm处的吸光度，根据Beer-Lambert定律计算蛋白质浓度。这种方法具有简便、灵敏、快速的优点，适合在小型装置中采用，并可以迅速得出结果。主要难点在于排除嘌呤、嘧啶、核酸等物质的干扰以及提高检测的准确度。

经过实验，我们发现在280-285nm区间上紫外光光吸光度随着蛋白质浓度变化有良好的相关性。

作为面向对象为病患的一款产品，我们的目标蛋白质浓度有效检测区间大约是0-4000毫克/升，经过实验论证，在该浓度下，蛋白质浓度对于尿液的吸光度有极为显著的影响，其他杂质的干扰较少，可以排除。

同时，相较于传统尿液分析仪和尿液有形成分数字分析仪，本设备成本低廉，能够大幅降低用户的检测费用。同时，设备可重复使用，减少了试纸等消耗品的使用成本，进一步提升了经济性。

\subsection*{3.3 技术路线设计}
我们产品的技术设计主要分为以下6部分：紫外光源、菲涅尔透镜、样品容器、光电二极管及放大电路、ESP32单片机、显示部分。

通过紫外LED光源发射散射的紫外光，再通过透镜的汇聚作用将LED产生的紫外光聚焦，起到聚焦光信号的作用，使得接收到的光强变化更为明显。

样品容器为石英比色皿，加入被测液体后可放置到指定位置。当样品吸收一定量的紫外光后，根据接受光信号的强弱可以判断样品中蛋白质浓度的情况，接收的光信号越强则样品中蛋白质浓度越低。

放大电路为设计的PCB板，主要功能是把光电二极管接收到的电流信号转化为电压信号并放大，使得单片机能够接收并处理该信号。
最后由ESP32单片机处理信号并处理，然后发送至电脑接收端并以网页形式展出，也可以在装置的显示屏上显示，显示信息包括蛋白质浓度与判定尿蛋白浓度区间，具体表示为“正常”、“+”、“++”等。

整个装置固定于一个体积较小的盒子内部，防止外部光线干扰检测的同时，可保障光路的稳定性，从而确保检测结果的准确性，同时也体现了便携以及装置小巧的特点。

\begin{figure}[H]
    \centering
    \fbox{\begin{minipage}{0.9\textwidth}
    \centering
    \vspace{0.5cm}
    \textcolor{gray}{\LARGE [此处放置图片: 技术流程图]}\\[0.5cm]
    \textbf{【图片内容】：} 描述从紫外光源到单片机处理，再到显示屏展示的技术数据流向图或光路示意图。\\
    \textbf{【图片作用】：} 结构化地展示了产品的工作原理和信号处理的先后顺序，使技术路线更加清晰易懂。\\
    \vspace{0.5cm}
    \end{minipage}}
    \caption{技术流程图}
\end{figure}

\begin{figure}[H]
    \centering
    \fbox{\begin{minipage}{0.9\textwidth}
    \centering
    \vspace{0.5cm}
    \textcolor{gray}{\LARGE [此处放置图片: 网页示意图]}\\[0.5cm]
    \textbf{【图片内容】：} 展示设备将数据传输至电脑端后，以网页形式呈现的检测结果界面。\\
    \textbf{【图片作用】：} 展示了产品的软件交互界面，证明产品具备上位机数据呈现和远程监控的能力。\\
    \vspace{0.5cm}
    \end{minipage}}
    \caption{网页示意图}
\end{figure}

\subsection*{3.4 产品功能}
本便携式尿液中蛋白质检测设备是一种利用紫外分光光度法，通过280-285nm波长的紫外光透射率测定尿液中的蛋白质浓度。可以实现尿液样品的蛋白质浓度分析并通过显示屏或电脑网页等方式向用户提交分析结果。

本设备体积小巧，重量轻，便于携带，非常适合家庭、社区医疗和移动医疗服务。用户只需将尿液样本放入设备，启动检测程序，几分钟后即可获得结果，无需专业操作技能，极大地方便了用户使用。

\begin{figure}[H]
    \centering
    \fbox{\begin{minipage}{0.9\textwidth}
    \centering
    \vspace{0.5cm}
    \textcolor{gray}{\LARGE [此处放置图片: 检测电压与蛋白质关系图]}\\[0.5cm]
    \textbf{【图片内容】：} 一个折线图或散点图，横坐标为蛋白质浓度，纵坐标为检测到的电压值，展示两者之间的关系。\\
    \textbf{【图片作用】：} 用实验数据证明了通过电压测量蛋白质浓度的科学依据，展示了测量结果的线性相关性和设备的可靠性。\\
    \vspace{0.5cm}
    \end{minipage}}
    \caption{检测电压与蛋白质关系}
\end{figure}

\begin{table}[H]
    \centering
    \caption{蛋白质浓度与对应指标数值}
    \renewcommand{\arraystretch}{1.2}
    \begin{tabular}{|c|c|c|c|c|}
    \hline
    蛋白质浓度（mg/L） & 0-100 & 100-670 & 670-2300 & >2300 \\
    \hline
    对应尿检指标 & 正常 & + & ++ & +++ \\
    \hline
    \end{tabular}
\end{table}

由上述实测图表可知，产品在尿检对应的蛋白质浓度上的电压检测结果呈现出较好的线性关系，能够提供给用户可靠的测试数据。

综上所述，本产品是一款能够较精准测量尿液中蛋白质浓度的产品，通过紫外分光光度的原理达到快速精确响应的效果，并且产品制造成本低，制作简便，因此具有较高的商业价值。

\section*{四、 财务分析}
\subsection*{4.1 成本核算}
设备生产成本包括紫外LED、菲涅尔透镜、比色皿、光电二极管、PCB、单片机、外壳，明细参考下表：

\begin{table}[H]
    \centering
    \caption{各元件成本汇总}
    \renewcommand{\arraystretch}{1.2}
    \begin{tabular}{|p{8cm}|c|}
    \hline
    \textbf{所需元件} & \textbf{价格（单位：元）} \\
    \hline
    紫外 LED（深圳泰溢光电，280-285nm） & 12 \\
    \hline
    菲涅尔透镜（直径20mm，焦距15mm） & 11.88 \\
    \hline
    石英比色皿 & 8.5 \\
    \hline
    光电二极管（杭州柠光科技，200-1100nm） & 57 \\
    \hline
    PCB & 3 \\
    \hline
    ESP32 单片机 & 6 \\
    \hline
    外壳（3D打印） & <1 \\
    \hline
    \textbf{总计} & \textbf{$\approx$ 99} \\
    \hline
    \end{tabular}
\end{table}

\subsection*{4.2 财务需求及预测}
大量资金在研发和生产初期需要被投入到设备购置、原材料采购等方面。同时，还需要预留市场营销费用，用于品牌推广和市场开拓。为了确保日常运营顺利，还需准备一定的营运资金，用于日常运营维护开支。

收入预测方面，可以根据市场需求预测每年的销售数量和单价，预计从销售收入中获得主要收入来源。还可以通过提供设备维护和升级服务获得售后服务收入，以及与医疗机构和养老院等合作项目的收入。

在成本预测方面，随着生产规模的扩大，单位制造成本预计会逐渐下降，而运营成本（包括人力、营销和行政等）可能会逐年增长。通过计算净利润，可以预测企业在首年可能处于亏损状态，但随着市场的拓展，预计第二年能接近收支平衡，并在第三年开始实现盈利。

\subsection*{4.3 风险预测}
从市场角度，风险包括市场接受度、竞争风险和政策风险。需评估产品能否被市场广泛接受，消费者是否愿意购买，并密切关注市场上是否有更具竞争力的产品，以制定相应的竞争策略。此外，政府对医疗器械行业的政策变化也可能影响销售。需确保产品技术的领先性和使用中的稳定性与准确性，符合国家设定的医疗标准。其他风险则包括知识产权纠纷和合规风险，需确保产品技术和设计不涉及知识产权问题，并符合相关法律法规和行业标准。

\section*{五、 营销手段及销售策略}
\subsection*{5.1 市场定位}
便携式尿液蛋白质检测设备以其高效、精准、经济的特点，主要定位于家庭健康管理、养老院和社区医疗机构等多个领域。

\textbf{（1）家庭健康管理}\\
随着人们健康意识的提高，家庭健康管理需求日益增长。便携式尿液蛋白质检测设备特别适用于家庭用户的日常健康监测需求。对于糖尿病、高血压和慢性肾病等慢性病患者，尿蛋白水平是重要的健康指标。该设备能帮助患者在家中随时监测尿液中的蛋白质含量，及早发现病情变化，减少突发疾病的风险。并且本设备可使患者无需频繁前往医院进行检查，减少了时间和精力的浪费。

\textbf{（2）养老院和社区医疗机构}\\
养老院和社区医疗机构面临着大量老年人和慢性病患者的健康监测需求。便携式尿液蛋白质检测设备的引入，可以显著提高这些机构的健康监测效率，并减少医疗成本。
同时，该设备能立刻提供检测结果，使医护人员能够更快速地对老年人的健康状况做出判断和反应，提高工作效率。并且由于该设备操作简便，无需复杂的专业培训，养老院和社区医疗机构可以减少对外部实验室检测的依赖，节省大量检测费用。同时，该设备的可重复使用性进一步降低了单次检测的成本。

\subsection*{5.2 营销手段}
\textbf{（1）官网营销}\\
建立便携式尿液蛋白质检测设备的官方网站，作为主要的营销和信息发布平台。官网将包含以下内容：设备的功能、使用方法、技术参数和价格信息、提供详细的使用说明和常见问题解答，帮助用户更好地理解和使用设备、展示真实用户的评价和使用案例，增强产品的可信度和吸引力、定期发布有关健康监测、慢性病管理等相关的文章和新闻，提升网站的专业性和权威性。

\textbf{（2）自媒体营销}\\
在抖音、快手、B站等短视频网站上发布便携式尿液蛋白质检测设备的相关视频，通过视频传播科学健康监测的效果，并提升品牌口碑。此外，邀请医疗专家、健康博主等专业人士参与内容创作，制作科学性强、易于理解的短视频，介绍设备的功能和使用方法。拍摄真实用户的体验视频，展示设备在日常健康监测中的实际效果，增强视频的真实性和吸引力。

\subsection*{5.3 品牌推广}
讲述品牌的创立背景和发展历程，展示企业的使命和愿景，增强品牌的温度和亲和力。通过品牌故事的传播，使用户对品牌产生情感共鸣和信任感。还可以与健康管理、慢性病管理等相关行业的知名企业和机构进行跨界合作，联合开展健康讲座、义诊活动等，提升品牌的影响力和美誉度。与此同时，积极参与社会公益活动，如免费健康检测、社区义诊等，通过实际行动展示企业的社会责任感，可以有效增强品牌的公信力和美誉度。

\section*{六、 团队情况}
孙以清，男，复旦大学电子信息科学与技术专业22级本科生。擅长3D建模和模型打印，主要负责产品的外观设计和内部结构设计。近年来多次参加各类竞赛，如全国大学生电子设计竞赛等。

周润，男，复旦大学电子信息科学与技术专业22级本科生。擅长Python与MATLAB编程，主要承担MCU搭建及上位机编程任务。曾参加全国大学生电子设计竞赛等比赛。

刘志杰，男，复旦大学电子信息科学与技术专业22级本科生。擅长模拟电子线路设计、搭建、测试，主要承担产品中电路设计部分。作为主力队员参加各类竞赛，如全国大学生数学建模竞赛、全国大学生电子设计竞赛等。

\section*{七、 总结及展望}
\subsection*{7.1 总结}
本便携式尿液中蛋白质检测设备是一种利用紫外分光光度法，通过280-285nm波长的紫外光透射率测定尿液中的蛋白质浓度。可以实现尿液样品的蛋白质浓度分析并通过显示屏或电脑网页等方式向用户提交分析结果。本产品有简便、成本低廉且可重复使用的优点，可以为用户提供快速、经济、高效的尿液蛋白质含量检测方案，具有较高商业价值以及良好的市场前景。

\subsection*{7.2 产品改进及展望}
本产品作为尿液检测设备，可以在以下方面提升：通过多波长组合的方式测量尿液中蛋白质浓度，并进行算法处理，提供更精确的检测结果，使其他成分干扰进一步减小。此外，还可以用类似的方式对其他尿液成分进行分析，提供更全面的检测。我们还预期将产品智能化，如为其设计微信小程序，将患者的尿蛋白指标提供给患者家属，可以方便监测患者的健康状况，也作下记录，为治疗提供参考。

\section*{八、 参考文献}
[1]孙玉标.临床尿常规检验的方法进展及评价 [J]. 中国医药指南， 2014, 12(17), 69 -70.\\
\null[2]上海市肾内科临床质量控制中心专家组 . 慢性肾脏病早期筛查、诊断及防治指南 （2